% preamble.tex — shared preamble for the SU(2) mass-gap monograph.
%
% Loaded by formal/main.tex BEFORE \begin{document}.  The class itself
% lives in docs/arthur_tex_style/arthur-book.cls (Tufte-style book class
% with margin-citation support and section-numbered theorem env's; it
% already loads amsmath, amssymb, amsthm, mathtools, hyperref, and the
% theorem environments theorem/lemma/proposition/corollary/definition/
% example numbered by section).
%
% The Makefile's `paper` target sets TEXINPUTS to find the class without
% a system install.  Run scripts/install_arthur_tex_style.sh once to make
% the class available globally.

% ---------- Additional packages ----------
% (arthur-book.cls already loaded amsmath et al; only add what's new.)

\usepackage[utf8]{inputenc}
\usepackage{xcolor}
\usepackage{enumitem}
\usepackage{listings}

% Allow bare `_` in text mode (the project's bibliography `note` fields
% reference Python identifiers like `sigma_hat_trajectory.py`).  Must
% come BEFORE biblatex.
\usepackage{underscore}

\usepackage[backend=biber,style=numeric-comp,maxbibnames=99]{biblatex}
\addbibresource{references.bib}

% cleveref must load AFTER hyperref (loaded by the class).
\usepackage[capitalise,noabbrev]{cleveref}

% Graphics: pdflatex does not embed SVG directly.  The Makefile's
% `svg-to-pdf` target pre-converts every SVG under visualizations/
% to a sibling PDF via rsvg-convert; \includegraphics then resolves
% by stem (e.g. {casimir-spectrum} -> visualizations/pilot/su2/casimir-spectrum.pdf).
\graphicspath{%
  {../visualizations/pilot/su2/}%
  {../visualizations/structure/}%
  {../visualizations/transfer/}%
}

% ---------- Extra theorem-like environments ----------
% arthur-math.sty already provides theorem, lemma, proposition, corollary,
% definition, example (section-numbered).  Add domain-specific ones.

\newtheorem{conjecture}[theorem]{Conjecture}
\newtheorem{milestone}[theorem]{Milestone}
\newtheorem{remark}[theorem]{Remark}
\newtheorem{convention}[theorem]{Convention}
\theoremstyle{arthurdefinition}
\newtheorem*{notation}{Notation}

% ---------- Project macros ----------

% Group theory.
\newcommand{\SUtwo}{\mathrm{SU}(2)}
\newcommand{\SUthree}{\mathrm{SU}(3)}
\newcommand{\SUN}{\mathrm{SU}(N)}
\newcommand{\Casimir}{C_2}

% RG / measures.
\newcommand{\Omk}{\Omega_k}
\newcommand{\Omkpone}{\Omega_{k+1}}
\newcommand{\Omksf}{\Omega_k^{\mathrm{sf}}}
\newcommand{\Omklf}{\Omega_k^{\mathrm{lf}}}
\newcommand{\Omkcap}{\Omega_k^{\cap}}
\newcommand{\muk}{\mu_k}
\newcommand{\muinf}{\mu_\infty}
\newcommand{\Tb}{T_b}

% Sheaves / KP.
\newcommand{\Ak}{\mathcal{A}_k}
\newcommand{\Sk}{\mathcal{S}_k}
\newcommand{\Pk}{\mathcal{P}_k}
\newcommand{\Ik}{\mathcal{I}_k}

% String tension / budget.
\newcommand{\sigk}{\sigma_k}
\newcommand{\sigh}{\widehat{\sigma}}
\newcommand{\sighk}{\widehat{\sigma}_k}
\newcommand{\sighinf}{\widehat{\sigma}_\infty}
\newcommand{\etak}{\eta_k}
\newcommand{\Dcap}{\Delta_{\cap}}
\newcommand{\Thetak}{\Theta_k}

% Cluster expansion.
\newcommand{\Ksf}{K_{\mathrm{sf}}}
\newcommand{\Klf}{K_{\mathrm{lf}}}
\newcommand{\Kcap}{K_{\cap}}
\newcommand{\rhok}{\rho_k}
\newcommand{\rhosf}{\rho_{\mathrm{sf}}}
\newcommand{\rholf}{\rho_{\mathrm{lf}}}
\newcommand{\rhocap}{\rho_{\cap}}

% Continuum / OS.
\newcommand{\Hphys}{\mathcal{H}_\infty}
\newcommand{\Hop}{H_{\mathrm{op}}}
\newcommand{\Eplus}{\mathcal{E}_+}
\newcommand{\osPair}{\langle\cdot,\cdot\rangle_{\mathrm{OS}}}
\newcommand{\massgap}{m_0}

% Constructive rationals.
\newcommand{\Qp}{\mathbb{Q}'}

% Fiducials (used in Appendix D and Ch. 3.1, 3.2, 3.4).
\newcommand{\clambda}{c_\lambda}
\newcommand{\muplaq}{\mu_{\mathrm{plaq}}}
\newcommand{\cpolymer}{c_{\mathrm{polymer}}}
\newcommand{\Amin}{A_{\min}}
\newcommand{\cchain}{c_{\mathrm{chain}}}

% Code references — keep tight typography for path names.
\newcommand{\codepath}[1]{\texttt{#1}}
\newcommand{\modref}[1]{\codepath{#1}}

% Inline "code-anchor" marker: a small bracketed tag in the margin or
% inline indicating the Python/Lean entry point of an equation or
% theorem.  Renders as "[code: path::name]" in slate-gray small font.
\newcommand{\codeanchor}[1]{%
  \hspace{0.5em}\textsuperscript{\textcolor{gray!70!black}{\footnotesize[\texttt{#1}]}}%
}

% ---------- Legacy / chapter-internal macros ----------
% These mirror the macros that used to live in each chapter file's own
% preamble.  Kept here so the chapter bodies can be \include'd verbatim
% (after stripping their per-file \documentclass...\begin{document}
% boilerplate).  Use \providecommand to be safe against name clashes.

% Calligraphic letter shortcuts.
\providecommand{\Acal}{\mathcal{A}}
\providecommand{\Bb}{\mathcal{B}}
\providecommand{\Ecal}{\mathcal{E}}
\providecommand{\Ecalp}{\mathcal{E}_+}
\providecommand{\Fcal}{\mathcal{F}}
\providecommand{\Hcal}{\mathcal{H}}
\providecommand{\Ical}{\mathcal{I}}
\providecommand{\Ncal}{\mathcal{N}}
\providecommand{\Pcal}{\mathcal{P}}
\providecommand{\Scal}{\mathcal{S}}
\providecommand{\Tcal}{\mathcal{T}}

% Non-negative reals and extended reals.
\providecommand{\Rnn}{\mathbb{R}_{\ge 0}}
\providecommand{\Rnnext}{\overline{\mathbb{R}}_{\ge 0}}
\providecommand{\Rd}{\mathbb{R}^d}
\providecommand{\Rfour}{\mathbb{R}^4}

% Configuration-space shortcuts (no clashes with the main macros above).
\providecommand{\Omsf}{\Omega_k^{\mathrm{sf}}}
\providecommand{\Omlf}{\Omega_k^{\mathrm{lf}}}
\providecommand{\Omcap}{\Omega_k^{\cap}}
\providecommand{\Omkp}{\Omega_{k+1}}
\providecommand{\Ominf}{\Omega_\infty}
\providecommand{\Omegavac}{\Omega_0}

% Block-spin pushforward star.
\providecommand{\Tbstar}{(\Tb)_{\!*}}

% Cylinder set and edge set.
\providecommand{\Cyl}{\mathrm{Cyl}}
\providecommand{\Ek}{E_k}

% OS abbreviation.
\providecommand{\OS}{\mathrm{OS}}

% Inner-product shortcuts.
\providecommand{\Iprod}[2]{\langle #1, #2 \rangle}
\providecommand{\IprodOS}[2]{\langle #1, #2 \rangle_{\mathrm{OS}}}

% Time-evolution operators (used in os-reconstruction).
\providecommand{\Tt}{T(t)}
\providecommand{\Ut}{U(t)}
